Generating \gls{DNP} group parameter data in a reactor-agnostic manner requires a reactor-agnostic generation process.
Fully detailing the flow regime with a large number of parameters is computationally infeasible.
Instead, the flow is simplified to two parameters: an in-core average residence time, $\tau_{in}$, and an ex-core residence time, $\tau_{ex}$.

This approach captures fuel salt advection around the primary loop and approximates turbulent and diffusive mixing.
The terms derive from the averaged behavior of a fuel salt control volume, treating thermal-hydraulic effects, such as vortices, as changes in residence times.
This method allows for approximate fitting of any circulating fuel reactor, requiring only the generation of relevant parameters for determining the associated \gls{DNP} group parameters.

The other required parameters describe chemical removal, which varies by fission product and reprocessing rate.
Refueling is not considered in this work, as its effects are negligible over the short time frames analyzed.
An irradiation time of approximately seven minutes suffices to reach equilibrium \glspl{DNP} concentrations, as the saturation time frame is on the order of multiple half-lives of the longest-lived \gls{DNP}.
Elements with cycle times on the order of months to years exhibit negligible changes, as discussed in Section \cite{sec:large-tcyc}.
Elements with changes on the order of seconds to days contribute sufficiently and must be included for accurate data generation.

The \gls{MSBR} reprocessing scheme serves as the basis for cycle times in this work, as shown in Table \ref{tab:msbr_cycle_times}.
Rare earths, semi-noble metals, volatile fluorides, discard, and higher nuclides have negligible effects over the time scales studied.
Section \ref{sec:large-tcyc} examines these longer cycle times to confirm their negligible impact.
Section \ref{sec:repr-scale} scales the cycle times to analyze how reprocessing rates influence the group parameters.

\begin{table}[H]
\renewcommand{\arraystretch}{1.25}
\caption{MSBR Online Reprocessing Cycle Times from \cite{robertson_conceptual_1971}.}
\label{tab:msbr_cycle_times}
\begin{center}
\begin{tabular}{ p{0.27\linewidth}  p{0.25\linewidth}  p{0.35\linewidth} }
 \hline
 Reprocessing Group & Element(s) & Cycle Time (Full Power)\\
 \hline
 Rare Earths & Y, La, Ce, Pr, Nd, Pm, Sm, Gd & 50 days\\
 Rare Earths & Eu & 500 days\\
 Noble Metals  & Se, Nb, Mo, Tc, Ru, Rh, Pd, Ag, Sb, Te & 20 seconds\\
 Seminoble Metals & Zr, Cd, In, Sn &  200 days\\
 Gases & Kr, Xe & 20 seconds\\
 Volatile Fluorides & Br, I & 60 days\\
 Discard & Rb, Sr, Cs, Ba & 3435 days\\
 Salt Discard & Th, Li, Be, F & 3435 days\\
 Protactinium & Pa & 3 days\\
 Higher Nuclides & Np, Pu & 16 years\\
 \hline
\end{tabular}
\end{center}
\end{table}

The parameters used include incident neutron energy, target fissile nuclide, in-core residence time, ex-core residence time, reprocessing scheme, and reprocessing scaling vector.
Each unique parameter adds a dimension to the data, exponentially increasing the computational cost of generating group parameters for every combination.
This work uses the \gls{MSBR} reprocessing scheme with a single scaling value applied to each reprocessing group instead of a unique scaling value for each.
The incident neutron energies that irradiate the sample are 0.0253 $eV$ and 1 $MeV$.
Section \ref{sec:res-time-var} determines the minimum number of residence times and their values required to capture general trends.
